% Created 2026-09-30 Wed 09:43
% Intended LaTeX compiler: pdflatex
\documentclass[11pt]{article}
\usepackage[utf8]{inputenc}
\usepackage[T1]{fontenc}
\usepackage{graphicx}
\usepackage{longtable}
\usepackage{wrapfig}
\usepackage{rotating}
\usepackage[normalem]{ulem}
\usepackage{amsmath}
\usepackage{amssymb}
\usepackage{capt-of}
\usepackage{hyperref}
\author{fcalado}
\date{\today}
\title{}
\hypersetup{
 pdfauthor={fcalado},
 pdftitle={},
 pdfkeywords={},
 pdfsubject={},
 pdfcreator={Emacs 29.3 (Org mode 9.6.15)}, 
 pdflang={English}}
\begin{document}

\tableofcontents

\section{week 6: Conference paper workshop day I}
\label{sec:orgb5a675b}

Agenda
\begin{itemize}
\item work on papers -> toward annotated bibliography
\item check in with me
\item check in with larger class
\item goal for today: leave class with 15-20 papers inside a zotero folder
\end{itemize}

\subsection{introducing annotated bibliography assignment - 10 minutes}
\label{sec:org8aaf9ef}
what's next with the papers: 
\begin{itemize}
\item 5 sources per person, 4 at the least (syllabus says 10-15, but was
for 2-3 person groups).
\item helping to get the ground research
\begin{itemize}
\item to form the perimeters of what's already been done in your topic
\item help you to further refine your questions
\item if you're doing a review type of paper, will also form your
analysis as well
\end{itemize}
\item will also form the basis of the lit review for the paper
\item needs to include:
\begin{itemize}
\item a short summary of the source (3 sentences), and
\item an explanation of how it may be useful for the project (1-2
sentences).
\end{itemize}
\item also due, in two weeks, a detailed outline of the paper.
\end{itemize}

\subsection{instructions for annotated bibs - 10 minutes}
\label{sec:org9dcec5c}
\begin{itemize}
\item how to read a conference paper: 3 passes.
\begin{itemize}
\item first pass:
\begin{itemize}
\item read only the most important sections, 10 - 15 minutes per paper
\item abstract, introduction, headings and subheadings, conclusions.
\item ask 2 things:
\begin{itemize}
\item category: what type of paper is this? a review of work? an
evaluation? an experiment? an analysis of an existing system?
\item contribution: what are the paper’s main contributions?
\end{itemize}
\end{itemize}
\item second pass:
\begin{itemize}
\item read the paper in full, up to 1 hour per paper.
\item where you will identify any potential gaps or areas of further research.
\item how is this useful for your project?
\end{itemize}
\item third pass:
\begin{itemize}
\item re-create the paper at each step, question every assumption.
\end{itemize}
\end{itemize}
\item places to look:
\begin{itemize}
\item FAccT proceedings
\item ACM digital library
\item ACL anthology
\item google scholar
\end{itemize}
\item strategies for looking:
\begin{itemize}
\item checking keywords, conferences, topics.
\end{itemize}
\end{itemize}

\subsection{finding articles activity - 30 - 45 minutes}
\label{sec:orgb4ce6a2}
\begin{itemize}
\item goal is to leave class with a list of 15-20 sources, with 4-5
assigned to each member in your group.

\item assign each person a turf:
\item by sub topic or by venue (conference)

\item each person completes a first pass on a number of articles on their
turf
\item final goal is to have at least 4-5 articles
\end{itemize}

\subsection{break}
\label{sec:org4418d7d}
\subsection{group check-ins - 45 minutes - 1 hour}
\label{sec:org827d097}
First, think about:
\begin{itemize}
\item what is the priority for your group at this point in the process
\begin{itemize}
\item what do you need to discuss with your group today
\end{itemize}
\end{itemize}
Then, in groups:
\begin{itemize}
\item decide on regular meeting schedule (if you don't have one already)
\item continue to work on finding papers
\end{itemize}

\subsection{big group check in}
\label{sec:orgf2a19cb}
\begin{itemize}
\item what is your area of research?
\end{itemize}
\end{document}
